\subsection{Motivation}

We performed empirical studies to present our motivations. The environmental details of all motivation experiments will be elaborated upon in §4.1.\par


\subsubsection{Sparse, heterogeneous, and dynamic metrics}
\begin{figure}[t]
\centering
\small
\includegraphics[width=\linewidth]{image/c2_sparsity_bars.pdf}
\caption{Sparsity of service-level metrics.} \label{m1.1.pdf}
\end{figure}

\begin{figure}[t]
\centering
\small
\includegraphics[width=\linewidth]{image/c2_latency_availability.pdf}
\caption{Availability of p90 latency (black=present, white=missing).} \label{m1.2.pdf}
\end{figure}

\begin{figure}[t]
\centering
\small
\includegraphics[width=\linewidth]{image/c2b_service_heterogeneity.pdf}
\caption{Heterogeneity among services in terms of CPU, memory, and execution time.} \label{m1.3.pdf}
\end{figure}

\begin{figure}[t]
\centering
\small
\includegraphics[width=\linewidth]{image/c2r_vertex_sparsity_variance.pdf}
\caption{Per-service invocation sparsity (left) and exec-time variance (right).} \label{m1.4.pdf}
\end{figure}

Currently, root cause localization methods typically rely on monitoring metrics to characterize service states and distinguish anomalous behaviors from fluctuations under a normal workload. However, agent services exhibit highly sparse, heterogeneous, and dynamic monitoring patterns, making it difficult to extract reliable failure features. As shown in Fig. \ref{m1.1.pdf} and Fig. \ref{m1.2.pdf}, even under a 5-user normal workload, approximately 79\% of latency metrics are missing and approximately 79\% of QPS metrics are zero. Service observability is also highly uneven, with the entry planner appearing in approximately 86\% of requests, while most tool agents are invoked in fewer than 25\%. Meanwhile, as shown in Fig. \ref{m1.3.pdf}, CPU utilization, execution time, and memory usage vary substantially across services, with execution time differing by approximately 55$\times$ and CPU utilization by approximately 3$\times$. Even within individual services, metrics under a normal workload exhibit large intrinsic fluctuations. The CV of network and latency-related metrics can reach 1--4, while the execution-time CV reaches approximately 2.8 for OCR and 0.7 for the planner (Fig. \ref{m1.4.pdf}). These characteristics make conventional statistical thresholds and common metric scales ineffective for distinguishing failures from variations under a normal workload, particularly for low-frequency services whose anomalous behaviors can be masked by sparse and inactive metrics. Consequently, accurately identifying anomalous services and extracting discriminative failure features from agent-service monitoring data remains challenging.


\begin{mybox}
\textbf{Insight 1:} It is essential to develop a service-aware and dynamic metric modeling mechanism that can identify discriminative failure features from sparse, heterogeneous, and highly variable metrics.
\end{mybox}



%统一论点：在面向 LLM 多智能体系统（以文档处理型 agent-network 为测试床）的可观测性数据上，我们采集了指标(metrics)、调用链(trace)、日志(log)三个维度，覆盖 1/3/5 并发用户与单/多副本多种部署架构。分析表明指标滞后性使异常时间窗口难以确定，进一步增大了Agent Service的根因定位难度。

\subsubsection{Dynamic metric lag and anomaly time windows}

\begin{figure}[t]
\centering
\small
\includegraphics[width=\linewidth]{image/recheck_lag_ccf.pdf}
\caption{Cross-correlation analysis of QPS-driven resource utilization, revealing metric lag in CPU and memory metrics.} \label{m2.1.pdf}
\end{figure}

%TODO下面两个图还是要重新搞user=10的新数据
\begin{figure}[t]
\centering
\small
\includegraphics[width=\linewidth]{image/c1r_completion_by_depth.pdf}
\caption{Completion time accumulates and disperses with invocation chain depth.} \label{m2.2.pdf}
\end{figure}


\begin{figure}[t]
\centering
\small
\includegraphics[width=\linewidth]{image/c1r_chain_dynamism.pdf}
\caption{Distribution of chain length.} \label{m2.4.pdf}
\end{figure}

% Currently, existing root cause localization methods typically rely on predefined anomaly time windows to align anomalous metrics and identify potential root causes. However, monitoring metrics in agent services exhibit substantial and dynamic metric lag. As shown in Fig. \ref{m2.1.pdf}, under a 5-user normal workload, the peak cross-correlation between QPS and CPU indicates metric lag of 30s and 10s under single- and multi-replica deployments, respectively, while metric lag for memory reaches 35s and 60s. Under a 1-user normal workload, no apparent metric lag is observed because the startup-induced resource peak coincides with the initial workload peak. Thus, metric lag varies across different metrics, normal workload settings, and deployment configurations, making a fixed temporal offset insufficient for metric alignment.

% The temporal dependency is further amplified along dynamic agent invocation chains. As requests propagate through multiple agent services, downstream metrics exhibit metric lag of approximately 15--60s relative to the entry stage, with completion time and variance increasing with chain depth (Fig. \ref{m2.2.pdf}). Meanwhile, invocation chains span 1--9 levels and contain dozens of distinct execution paths, each with different service compositions and execution times (Fig. \ref{m2.4.pdf}, Table \ref{tab:m2.3}). Unlike traditional microservice systems with relatively stable service dependencies, agent services can dynamically change their execution paths even under the same normal workload. For example, the observable fraction of word-generation latency varies from 6.6\% to 55.8\%, while that of OCR changes from 59.1\% to 1.7\%, primarily due to different task combinations and execution paths.

Existing root cause localization methods typically rely on predefined anomaly time windows to align anomalous metrics and identify potential root causes. However, agent services exhibit substantial and dynamic metric lag. As shown in Fig. \ref{m2.1.pdf}, the lag between QPS and resource utilization varies across metrics, workloads, and deployment configurations: under a 5-user workload, CPU lag changes from 30s to 10s between single- and multi-replica deployments, while memory lag changes from 35s to 60s. This variability indicates that metric lag cannot be captured by a fixed temporal offset.​Such metric lag is further shaped by the execution characteristics of agent services. Requests traverse dynamic invocation chains with different service compositions and execution times. As shown in Figs. \ref{m2.2.pdf} and \ref{m2.4.pdf}, completion time and its variance increase with chain depth, while invocation chains span 1--9 levels and follow dozens of distinct execution paths. Consequently, metrics from different services may correspond to different stages of request execution, introducing heterogeneous temporal offsets along invocation chains. ​Determining anomaly time windows becomes more difficult in the presence of failures. As shown in Fig. \ref{m2.5.pdf}, CPU metrics during failures can substantially overlap with those of the pre-failure workload in the same run, and this overlap varies widely across failure types, loads, and faulted services, making failure-induced changes difficult to distinguish from normal fluctuations. Moreover, Fig. \ref{m2.6.pdf} shows that metric responses do not necessarily coincide with failure injection: degradation may emerge substantially later or even precede failure onset due to normal workload fluctuations. Failure propagation along invocation chains therefore introduces additional and non-uniform metric lag, making neither failure onset nor a fixed temporal offset a reliable boundary for anomaly windows.​These observations show that accurately estimating metric lag is a prerequisite for determining anomaly time windows and aligning anomalous metrics in agent services. Without accounting for metric lag introduced by service execution characteristics and failure propagation, fixed or uniformly shifted windows can incorporate unrelated fluctuations or miss failure-related metrics, undermining subsequent anomaly detection and root-cause localization.

\begin{figure*}[t]
\centering
\small
\includegraphics[width=\linewidth]{image/c2c_abn_vs_normal_overlap.pdf}
\caption{Variability and overlap of metrics under normal workload and failures. (a) Service-level metrics under different normal workloads, measured by CV. (b) CPU CV of the faulted service during each failure (bars) versus its CV under the pre-failure workload of the same run (ticks), for PDF-parsing (3 and 5 users) and for the planner and summarizer (5 users). (c) The percentage of CPU metrics during each failure that fall within the p10--p90 range of the pre-failure workload of the same run, for the same faulted services.} \label{m2.5.pdf}
\end{figure*}

% Determining anomaly time windows additionally requires distinguishing failure-induced changes in metrics from fluctuations under a normal workload. To examine this, we inject five representative failures (CPU stress, memory stress, network latency, pod failure, and pod kill) into the PDF-parsing service under a 3-user workload, and compare its metrics with those from the same service under the corresponding normal workload. As shown in Fig. \ref{m2.5.pdf}, CPU and memory metrics during failures often exhibit patterns similar to those under normal workloads. Such overlap makes it difficult to distinguish from normal fluctuations, making the agent services more susceptible to normal-variation noise when identifying anomalous time windows.
% CPU metrics fluctuate substantially under the normal workload and during failures: the CV is 1.69 under the normal workload and ranges from 0.46 to 3.36 across the five failure experiments. 
% all CPU and memory metrics during failure injection in the network-latency experiment fall within their respective p10--p90 ranges under the normal workload, as do 56.8\% of CPU metrics during pod kill. This overlap varies across metrics and failure types: during CPU stress, 18.9\% of CPU metrics but 100\% of memory metrics remain within the corresponding ranges under the normal workload. Thus, variability in metrics and departures from the normal workload range do not provide consistent temporal boundaries across failure types.


\begin{figure*}[t]
\centering
\small
\includegraphics[width=\linewidth]{image/c1_abn_window_timeline.pdf}
\caption{Metric lag between failure injection and changes in service metrics. (a)–(b) CPU utilization and success rates during CPU stress. (c) First success-rate reductions at the entry planner and downstream summarizer across failure experiments. Shading indicates the 181-s failure-injection period, and annotations indicate the timing of metric reductions relative to the end of failure injection.} \label{m2.6.pdf}
\end{figure*}


% As shown in Fig. \ref{m2.6.pdf}, changes in metrics do not necessarily coincide with failure injection. Metric degradation may occur substantially later than the injected failure, after other metrics have returned to their normal ranges, or may even precede failure onset due to normal workload fluctuations. Such metric lag varies across failure types and services along the invocation chain, making neither failure onset nor a fixed temporal offset a reliable boundary for anomaly windows. Therefore, accurately identifying anomalous metrics and root causes in agent services requires accounting for metric lag and failure propagation rather than relying on fixed or uniformly shifted time windows.

% Even when deviations in resource metrics are visible, their timing can differ substantially from changes in metrics elsewhere along the invocation chain. Across all five experiments, all available success-rate metrics of the service with failure equal 1.0; however, these metrics are entirely missing during CPU and memory stress. In contrast, success-rate degradation is observed at the downstream summarizer and, in four experiments, the entry planner.
% in the CPU-stress experiment, failure injection lasts 181s and CPU utilization peaks at 568mC, whereas the first success-rate reductions at the planner and summarizer appear 400s after failure onset, or 219s after the end of failure injection, when CPU utilization has already returned to the normal workload range. For CPU stress, memory stress, pod failure, and pod kill, the earliest reductions at these chain services occur 194--429s after the end of failure injection. The network-latency experiment instead contains a reduction 75s before failure onset and another 120s after failure onset, so its observed onset cannot be attributed unambiguously to the injected failure. These results show that anomaly time windows beginning at the first degradation in service metrics can miss failure injection, while shifting anomaly time windows backward by a common offset cannot accommodate differences in metric lag.

\begin{mybox}
\textbf{Insight 2:} Accurately estimating metric lag arising from the execution characteristics of agent services and failure propagation along invocation chains is essential for determining anomaly time windows, enabling reliable anomaly detection and subsequent root-cause localization.
\end{mybox}
